\section{The Coherence Term}\label{sec:coherence}

\subsection{Does the Curve Already Handle Incoherence?}

At first glance, the $a_k$ curve appears to handle the coherence problem from Section~\ref{sec:coherence-problem} on its own.
A policy that outputs pure noise has $a_k \approx \mathrm{const}$ for all $k$ (random tokens from one response do not help predict random tokens from another), so the curve stays high and $a_\infty$ is high but $C$ will be near zero.
The noise policy correctly gets a low $D_{Ca_n}$ because of the coherence term, not because of the curve alone.

A second case requiring $C$ is a \textbf{policy that has multiple distinct but individually incoherent modes}.
Consider a badly trained or modified model that produces several recognizably different types of low-quality output.
The base model $\theta$ can pick up on which type is being produced from context, but each type is still implausible text.
Without a separate plausibility term, the policy could score as meaningfully diverse.
A separate coherence term lets $D_{Ca_n}$ suppress this case as well.

\subsection{Definition}

The per-byte cross-entropy $h_\theta(r \mid p)$ from Eq.~\eqref{eq:pertok} can be converted to a geometric mean per-byte probability:
\begin{equation}
    2^{-h_\theta(r \mid p)} = \left(\prod_{t=1}^{|r|_{\mathrm{tok}}} \theta(r^t \mid r^{<t}, p)\right)^{1/\|r\|}
\end{equation}

This is the average per-byte probability the base model assigns to the response, living in $(0,1)$.

We define \textbf{coherence} as the geometric mean of this quantity across responses:
\begin{equation}\label{eq:coherence}
    C = 2^{-\frac{1}{n}\sum_{i=1}^{n} h_\theta(r_i \mid p)}
\end{equation}

This is interpretable as ``how likely, per byte, does the base model find this policy's outputs?''
Equivalently, $C$ is the reciprocal of the geometric-mean per-byte perplexity:\footnote{For scale, the bits/byte score is the average per-byte cross-entropy $\bar{\ell}$, related to coherence by $C = 2^{-\bar{\ell}}$. GPT-3 davinci on The Pile \citep{gao2020pile} scored $\sim$0.72 bits/byte and GPT-2 XL $\sim$1.05 bits/byte, giving $C \approx 0.61$ and $C \approx 0.49$ respectively; stronger base models reach lower bits/byte and correspondingly higher $C$.}
\begin{equation}
    C = \frac{1}{\mathrm{PPL}_\theta(\pi, p)}, \qquad \mathrm{PPL}_\theta(\pi, p) = 2^{\frac{1}{n}\sum_{i=1}^{n} h_\theta(r_i \mid p)}
\end{equation}

We chose the geometric mean (equivalently, $1/n$ inside the exponent) over the arithmetic mean $C_{\mathrm{AM}} = \frac{1}{n}\sum_i 2^{-h_\theta(r_i \mid p)}$.
By Jensen's inequality (since $2^{-x}$ is convex), $C_{\mathrm{AM}} \geq C$ with equality iff all $h_\theta(r_i \mid p)$ are identical.
The gap is precisely what makes the arithmetic mean problematic: if 9 out of 10 responses are incoherent and one is highly plausible, $C_{\mathrm{AM}}$ remains nonzero while $C$ is driven toward zero by the geometric averaging.

\paragraph{Distributional distance.} In expectation over $r \sim \pi$:
\begin{equation}
    \E_\pi\bigl[h_\theta(r \mid p)\bigr] = H_\pi(\cdot \mid p) + \KL(\pi \| \theta \mid p)
\end{equation}

So coherence is penalized both by genuine incoherence and by distributional distance from the base model.
For the purpose of evaluating modified policies, this is a feature: a policy that assigns probability mass to regions $\theta$ finds implausible is penalized accordingly, reflecting that the base model's probability landscape is treated as the reference.

\paragraph{Note on comparability.} Because all quantities are normalized per byte, the per-byte coherence $C$ is invariant to tokenizer choice: two base models $\theta_1, \theta_2$ with different tokenizers but identical predictive distributions will yield identical $C$ values.
What remains incomparable across different $\theta$'s is model quality: a better model assigns lower cross-entropy regardless of normalization.
This is by design: $\theta$'s predictive quality is the lens through which diversity is measured.

